% ===========================================================================
% PHỤ LỤC
% ===========================================================================

\chapter{CẤU TRÚC MÃ NGUỒN VÀ HƯỚNG DẪN TÁI LẬP KẾT QUẢ}
% Tiêu đề chạy rút gọn: tiêu đề chương đầy đủ dài hơn một dòng đầu trang,
% khiến fancyhdr phải nới \headheight và làm lệch khung chữ của cả chương.
\markboth{Phụ lục A. Cấu trúc mã nguồn và hướng dẫn tái lập kết quả}{}
\label{app:source}

\section{Tổ chức mã nguồn}

Toàn bộ quy trình được hiện thực bằng Python và tổ chức thành một gói phần mềm có
cấu trúc phản ánh đúng mạch nội dung nghiên cứu của đề án. Mỗi gói con tương ứng
một mục của Chương~\ref{chap:xay-dung-he-thong}, nhờ đó người đọc có thể đối
chiếu trực tiếp giữa phần trình bày phương pháp và phần hiện thực.

\begin{verbatim}
config/config.yaml           Cấu hình tập trung toàn bộ quy trình
src/creditreview/
  config.py                  Nạp và truy cập cấu hình
  data/                      2.2 Thu thập và tổng hợp dữ liệu
    generator.py           Sinh dữ liệu mô phỏng (biến tiềm ẩn)
    reference.py           Bảng tham chiếu nghiệp vụ
    profiling.py           Thống kê mô tả, kiểm tra hợp lệ, tiêm lỗi
  preprocessing/             2.3 Tiền xử lý và làm sạch dữ liệu
    cleaning.py              Xử lý giá trị khuyết thiếu và ngoại lệ
    encoding.py              Chuẩn hóa và mã hóa biến
    splitting.py             Phân chia tập dữ liệu
  features/                  2.4 Trích xuất và lựa chọn đặc trưng
    definitions.py           Đăng ký đặc trưng theo nhóm nguồn dữ liệu
    extraction.py            Trích xuất đặc trưng đúng theo thời điểm
    selection.py             Lựa chọn tập đặc trưng qua ba bộ lọc
  labeling/                  2.5 Xây dựng nhãn và quy tắc phê duyệt
    target.py                Gán nhãn rủi ro, phân tích vintage
    cutoff.py                Xác định điểm cắt phê duyệt
  models/                    2.6, 2.7 Thiết kế và huấn luyện mô hình
    zoo.py                   Nhà máy mô hình (12 mô hình)
    train.py                 Huấn luyện, hiệu chuẩn, đánh giá
    tuning.py                Tinh chỉnh siêu tham số
    metrics.py               Bộ chỉ tiêu đánh giá
    woe.py                   Biến đổi WOE và thẻ điểm
    explain.py               Giải thích mô hình (SHAP, phép hoán vị)
    plots.py                 Sinh đồ thị
  decision/                  Bộ máy hỗ trợ ra quyết định
    rules.py                 Bộ luật nghiệp vụ
    policy.py                Chính sách hạn mức và định giá
    engine.py                Tổng hợp khuyến nghị theo thứ tự ưu tiên
  serving/                   Dịch vụ REST
    api.py                   Các điểm cuối nghiệp vụ và vận hành
    repository.py            Kho đặc trưng phục vụ tra cứu
  monitoring/
    drift.py                 Giám sát PSI và chất lượng theo thời gian
  dashboard/
    app.py                   Giao diện trực quan hóa kết quả
scripts/                     Bảy bước tái lập kết quả
results/                     Mô hình, bảng và đồ thị kết quả
thesis/                      Mã nguồn LaTeX của đề án
\end{verbatim}

\section{Hướng dẫn tái lập kết quả}

\subsection{Cài đặt môi trường}

\begin{verbatim}
pip install -r requirements.txt
\end{verbatim}

Yêu cầu Python phiên bản 3.10 trở lên.

\subsection{Chạy toàn bộ quy trình}

Cách nhanh nhất là chạy toàn bộ bằng một lệnh:

\begin{verbatim}
python scripts/run_all.py
\end{verbatim}

Hoặc chạy tuần tự từng bước để theo dõi kết quả trung gian:

\begin{verbatim}
# Bước 1 - Chuẩn bị dữ liệu: sinh, kiểm tra chất lượng,
#          gán nhãn, phân chia tập
python scripts/01_prepare_data.py

# Bước 2 - Phân tích Information Value và lựa chọn đặc trưng
python scripts/02_build_features.py

# Bước 3 - Tinh chỉnh siêu tham số (tùy chọn)
python scripts/03_tune_hyperparams.py --trials 14

# Bước 4 - Huấn luyện và so sánh 12 mô hình
python scripts/04_train_models.py

# Bước 5 - Mô phỏng quy trình kiểm duyệt và đánh giá chung
python scripts/05_simulate_evaluate.py

# Bước 6 - Xuất bảng LaTeX cho đề án
python scripts/06_export_latex_tables.py

# Bước 7 - Xuất từ điển dữ liệu các bảng đầu vào (Excel, tùy chọn)
python scripts/07_export_data_dictionary.py
\end{verbatim}

Hạt giống ngẫu nhiên được cố định trong \code{config/config.yaml}
(\code{project.seed}), bảo đảm kết quả lặp lại được giữa các lần chạy.

\subsection{Chạy nhanh để kiểm tra}

Để kiểm tra quy trình hoạt động mà không mất nhiều thời gian:

\begin{verbatim}
python scripts/01_prepare_data.py --n-customers 3000 \
    --n-applications 4500
python scripts/02_build_features.py
python scripts/04_train_models.py --no-stacking --no-shap \
    --models logistic_regression lightgbm
python scripts/05_simulate_evaluate.py --n-applications 1000 --no-api
\end{verbatim}

\subsection{Khởi động dịch vụ và giao diện}

\begin{verbatim}
# Dịch vụ REST hỗ trợ kiểm duyệt
uvicorn creditreview.serving.api:app --port 8000
#   Tài liệu API tự sinh: http://localhost:8000/docs

# Giao diện trực quan hóa kết quả
streamlit run src/creditreview/dashboard/app.py
\end{verbatim}

\subsection{Biên dịch đề án}

\begin{verbatim}
cd thesis
xelatex main.tex && xelatex main.tex
\end{verbatim}

Cần dùng XeLaTeX do tài liệu sử dụng tiếng Việt và phông Times New Roman. Chạy
hai lần để cập nhật mục lục và tham chiếu chéo.

Số liệu trong đề án được sinh tự động từ kết quả thực nghiệm qua
\code{scripts/06\_export\_latex\_tables.py}: chạy lại thực nghiệm là toàn bộ bảng
biểu và các con số trích dẫn trong phần diễn giải tự cập nhật, không phải sửa
tay.

\chapter{DANH MỤC LUẬT NGHIỆP VỤ}
% Tiêu đề chạy rút gọn: tiêu đề chương đầy đủ dài hơn một dòng đầu trang,
% khiến fancyhdr phải nới \headheight và làm lệch khung chữ của cả chương.
\markboth{Phụ lục B. Danh mục luật nghiệp vụ}{}
\label{app:rules}

Bảng dưới đây liệt kê đầy đủ bộ luật nghiệp vụ được cài đặt trong mô hình. Ba
mức hành động gồm: \code{HARD\_REJECT} -- từ chối, có quyền phủ quyết cao nhất;
\code{REVIEW} -- buộc chuyển thẩm định thủ công; \code{WARNING} -- ghi nhận vào
danh sách lý do nhưng không thay đổi khuyến nghị.

\begingroup
\small
\setlength{\tabcolsep}{4.0pt}
\begin{longtable}{@{}L{1.436cm}L{3.638cm}L{2.393cm}L{2.400cm}L{4.308cm}@{}}
\caption{Danh mục đầy đủ luật nghiệp vụ được cài đặt trong mô hình}\label{tab:rule-catalogue}\\
\toprule
\textbf{Mã luật} & \textbf{Tên luật} & \textbf{Nhóm} & \textbf{Hành động} & \textbf{Diễn giải} \\
\midrule
\endfirsthead
\multicolumn{5}{@{}l}{\small\textit{(tiếp theo trang trước)}}\\
\toprule
\textbf{Mã luật} & \textbf{Tên luật} & \textbf{Nhóm} & \textbf{Hành động} & \textbf{Diễn giải} \\
\midrule
\endhead
\bottomrule
\multicolumn{5}{r@{}}{\small\textit{(xem tiếp trang sau)}}\\
\endfoot
\bottomrule
\endlastfoot
R001 & Tuổi nằm ngoài khoảng cho phép [20; 62] & Điều kiện cơ bản & HARD\_\allowbreak REJECT & Khách hàng phải trong độ tuổi từ 20 đến 62 tại thời điểm đề nghị vay \\
R002 & Thu nhập dưới mức tối thiểu 6 triệu VND & Điều kiện cơ bản & HARD\_\allowbreak REJECT & Thu nhập hàng tháng không đạt mức tối thiểu của sản phẩm \\
R003 & Thâm niên công việc dưới 6 tháng & Điều kiện cơ bản & REVIEW & Thâm niên công việc chưa đủ để đánh giá tính ổn định của thu nhập \\
R101 & Nhóm nợ xấu nhất 12 tháng vượt nhóm 2 & Lịch sử tín dụng & HARD\_\allowbreak REJECT & Khách hàng phát sinh nợ xấu trong 12 tháng gần nhất \\
R102 & Đang có nợ nhóm 3 trở lên & Lịch sử tín dụng & HARD\_\allowbreak REJECT & Khách hàng đang có khoản nợ được phân loại từ nhóm 3 (nợ dưới tiêu chuẩn) \\
R103 & Có nợ đã xử lý rủi ro / xóa nợ & Lịch sử tín dụng & HARD\_\allowbreak REJECT & Lịch sử có khoản nợ đã được xử lý rủi ro hoặc xóa nợ \\
R104 & Số ngày quá hạn tối đa 12 tháng vượt 30 ngày & Lịch sử tín dụng & REVIEW & Từng chậm trả vượt ngưỡng cho phép trong 12 tháng gần nhất \\
R105 & Quan hệ tín dụng với hơn 6 tổ chức & Lịch sử tín dụng & REVIEW & Số lượng tổ chức tín dụng đang có quan hệ ở mức cao, tiềm ẩn rủi ro vay đảo nợ \\
R106 & Truy vấn tín dụng 3 tháng vượt 5 lần & Lịch sử tín dụng & WARNING & Khách hàng tìm kiếm tín dụng dồn dập tại nhiều tổ chức \\
R107 & Tỷ lệ sử dụng hạn mức trên 95\% & Lịch sử tín dụng & WARNING & Hạn mức tín dụng hiện có đã được sử dụng gần hết \\
R201 & Tỷ lệ nợ trên thu nhập (DTI) vượt 60\% & Khả năng trả nợ & HARD\_\allowbreak REJECT & Tổng nghĩa vụ trả nợ vượt quá khả năng chi trả từ thu nhập \\
R202 & Thu nhập còn lại sau trả nợ không đủ chi phí sinh hoạt & Khả năng trả nợ & HARD\_\allowbreak REJECT & Sau khi trừ nghĩa vụ trả nợ và chi phí sinh hoạt tối thiểu, thu nhập còn lại âm \\
R203 & Tỷ lệ trả góp trên thu nhập (PTI) vượt 45\% & Khả năng trả nợ & REVIEW & Riêng khoản vay mới đã chiếm tỷ trọng lớn trong thu nhập \\
R204 & Số tiền vay vượt 60 lần thu nhập hàng tháng & Khả năng trả nợ & REVIEW & Quy mô khoản vay rất lớn so với thu nhập, cần thẩm định bổ sung \\
R301 & Tỷ lệ vay trên giá trị TSBĐ (LTV) vượt 85\% & Tài sản bảo đảm & REVIEW & Giá trị tài sản bảo đảm không đủ tỷ lệ bao phủ yêu cầu \\
R401 & Không lấy được dữ liệu CIC & Chất lượng dữ liệu & REVIEW & Hệ thống CIC không phản hồi - chuyển thẩm định thủ công theo quy trình dự phòng \\
R402 & Khách hàng thiếu lịch sử tín dụng (thin-file) & Chất lượng dữ liệu & WARNING & Không có/ít lịch sử tín dụng; đánh giá dựa chủ yếu vào dữ liệu hành vi giao dịch \\
R403 & Thiếu dữ liệu hành vi giao dịch & Chất lượng dữ liệu & REVIEW & Số giao dịch trong 180 ngày quá ít để tính đặc trưng hành vi đáng tin cậy \\
R501 & Dòng tiền vào thực tế thấp hơn nhiều so với thu nhập khai báo & Dấu hiệu bất thường & REVIEW & Khách hàng nhận lương qua tài khoản nhưng dòng tiền vào 30 ngày dưới 35\% thu nhập khai báo - cần xác minh thu nhập \\
R502 & Tỷ trọng chi tiêu nhóm rủi ro cao bất thường & Dấu hiệu bất thường & REVIEW & Trên 15\% dòng tiền ra thuộc nhóm giao dịch rủi ro cao \\
R503 & Không phát sinh giao dịch trong 60 ngày gần nhất & Dấu hiệu bất thường & WARNING & Tài khoản không hoạt động trong thời gian dài trước khi đề nghị vay \\
\end{longtable}
\endgroup


\chapter{DANH MỤC ĐẶC TRƯNG}
% Tiêu đề chạy rút gọn: tiêu đề chương đầy đủ dài hơn một dòng đầu trang,
% khiến fancyhdr phải nới \headheight và làm lệch khung chữ của cả chương.
\markboth{Phụ lục C. Danh mục đặc trưng}{}
\label{app:features}

Danh mục đầy đủ các đặc trưng được sử dụng, phân theo nhóm nguồn dữ liệu. Cột
``Nguồn dữ liệu'' cho biết đặc trưng được tính từ bảng dữ liệu nào.

\begingroup
\small
\setlength{\tabcolsep}{4.0pt}
\begin{longtable}{@{}L{3.761cm}L{2.121cm}L{2.314cm}L{1.061cm}L{4.918cm}@{}}
\caption{Danh mục đặc trưng đầu vào của mô hình}\label{tab:feature-catalogue}\\
\toprule
\textbf{Tên đặc trưng} & \textbf{Nhóm} & \textbf{Nguồn dữ liệu} & \textbf{Kiểu} & \textbf{Mô tả} \\
\midrule
\endfirsthead
\multicolumn{5}{@{}l}{\small\textit{(tiếp theo trang trước)}}\\
\toprule
\textbf{Tên đặc trưng} & \textbf{Nhóm} & \textbf{Nguồn dữ liệu} & \textbf{Kiểu} & \textbf{Mô tả} \\
\midrule
\endhead
\bottomrule
\multicolumn{5}{r@{}}{\small\textit{(xem tiếp trang sau)}}\\
\endfoot
\bottomrule
\endlastfoot
age & Thông tin cá nhân & HỒ SƠ KH & num & Tuổi khách hàng \\
gender & Thông tin cá nhân & HỒ SƠ KH & cat & Giới tính \\
marital\_\allowbreak status & Thông tin cá nhân & HỒ SƠ KH & cat & Tình trạng hôn nhân \\
num\_\allowbreak dependents & Thông tin cá nhân & HỒ SƠ KH & num & Số người phụ thuộc \\
education\_\allowbreak code & Thông tin cá nhân & HỒ SƠ KH & num & Trình độ học vấn (1-5) \\
occupation & Thông tin cá nhân & HỒ SƠ KH & cat & Nghề nghiệp \\
employer\_\allowbreak type & Thông tin cá nhân & HỒ SƠ KH & cat & Loại hình tổ chức làm việc \\
job\_\allowbreak tenure\_\allowbreak months & Thông tin cá nhân & HỒ SƠ KH & num & Thâm niên công việc (tháng) \\
monthly\_\allowbreak income & Thông tin cá nhân & HỒ SƠ KH & num & Thu nhập khai báo hàng tháng \\
province & Thông tin cá nhân & HỒ SƠ KH & cat & Tỉnh/thành cư trú \\
residence\_\allowbreak type & Thông tin cá nhân & HỒ SƠ KH & cat & Hình thức nhà ở \\
customer\_\allowbreak since\_\allowbreak months & Thông tin cá nhân & HỒ SƠ KH & num & Thời gian quan hệ với ngân hàng \\
has\_\allowbreak salary\_\allowbreak account & Thông tin cá nhân & HỒ SƠ KH & bin & Có tài khoản nhận lương \\
n\_accounts & Thông tin cá nhân & HỒ SƠ KH & num & Số tài khoản đang mở \\
total\_\allowbreak avg\_\allowbreak balance & Thông tin cá nhân & HỒ SƠ KH & num & Tổng số dư bình quân 3 tháng \\
product\_\allowbreak code & Tham số khoản vay & ĐƠN VAY & cat & Mã sản phẩm tín dụng \\
requested\_\allowbreak amount & Tham số khoản vay & ĐƠN VAY & num & Số tiền đề nghị vay \\
tenor\_\allowbreak months & Tham số khoản vay & ĐƠN VAY & num & Kỳ hạn vay (tháng) \\
interest\_\allowbreak rate & Tham số khoản vay & ĐƠN VAY & num & Lãi suất áp dụng (năm) \\
loan\_\allowbreak purpose & Tham số khoản vay & ĐƠN VAY & cat & Mục đích vay \\
channel & Tham số khoản vay & ĐƠN VAY & cat & Kênh tiếp nhận hồ sơ \\
has\_\allowbreak collateral & Tham số khoản vay & ĐƠN VAY & bin & Có tài sản bảo đảm \\
monthly\_\allowbreak installment & Tham số khoản vay & ĐƠN VAY & num & Khoản trả góp hàng tháng \\
ltv & Tham số khoản vay & ĐƠN VAY & num & Tỷ lệ vay trên giá trị TSBĐ \\
apply\_hour & Tham số khoản vay & ĐƠN VAY & num & Giờ nộp hồ sơ \\
apply\_\allowbreak dayofweek & Tham số khoản vay & ĐƠN VAY & num & Ngày trong tuần nộp hồ sơ \\
cic\_score & Lịch sử tín dụng (CIC) & BẢN TIN CIC & num & Điểm tín dụng CIC (300-850) \\
current\_\allowbreak group & Lịch sử tín dụng (CIC) & BẢN TIN CIC & num & Nhóm nợ hiện tại (1-5) \\
worst\_\allowbreak group\_\allowbreak 12m & Lịch sử tín dụng (CIC) & BẢN TIN CIC & num & Nhóm nợ xấu nhất 12 tháng \\
dpd\_\allowbreak max\_\allowbreak 12m & Lịch sử tín dụng (CIC) & BẢN TIN CIC & num & Số ngày quá hạn tối đa 12 tháng \\
num\_\allowbreak institutions & Lịch sử tín dụng (CIC) & BẢN TIN CIC & num & Số TCTD đang có quan hệ \\
num\_\allowbreak active\_\allowbreak contracts & Lịch sử tín dụng (CIC) & BẢN TIN CIC & num & Số hợp đồng còn hiệu lực \\
num\_\allowbreak closed\_\allowbreak contracts & Lịch sử tín dụng (CIC) & BẢN TIN CIC & num & Số hợp đồng đã tất toán \\
total\_\allowbreak outstanding & Lịch sử tín dụng (CIC) & BẢN TIN CIC & num & Tổng dư nợ tại các TCTD \\
credit\_\allowbreak utilization & Lịch sử tín dụng (CIC) & BẢN TIN CIC & num & Tỷ lệ sử dụng hạn mức \\
monthly\_\allowbreak debt\_\allowbreak payment & Lịch sử tín dụng (CIC) & BẢN TIN CIC & num & Nghĩa vụ trả nợ hàng tháng hiện hữu \\
num\_\allowbreak inquiries\_\allowbreak 3m & Lịch sử tín dụng (CIC) & BẢN TIN CIC & num & Số lần truy vấn tín dụng 3 tháng \\
has\_\allowbreak written\_\allowbreak off\_\allowbreak debt & Lịch sử tín dụng (CIC) & BẢN TIN CIC & bin & Có nợ đã xử lý rủi ro \\
oldest\_\allowbreak account\_\allowbreak months & Lịch sử tín dụng (CIC) & BẢN TIN CIC & num & Độ dài lịch sử tín dụng (tháng) \\
num\_\allowbreak bad\_\allowbreak debt\_\allowbreak institutions & Lịch sử tín dụng (CIC) & BẢN TIN CIC & num & Số TCTD từng phát sinh nợ xấu \\
is\_\allowbreak thin\_\allowbreak file & Lịch sử tín dụng (CIC) & BẢN TIN CIC & bin & Khách hàng thiếu lịch sử tín dụng \\
cic\_\allowbreak missing & Lịch sử tín dụng (CIC) & BẢN TIN CIC & bin & Không lấy được dữ liệu CIC \\
txn\_\allowbreak cnt\_\allowbreak \{W\}d & Hành vi giao dịch & GIAO DỊCH & num & Số giao dịch trong cửa sổ W ngày \\
inflow\_\allowbreak \{W\}d & Hành vi giao dịch & GIAO DỊCH & num & Tổng dòng tiền vào trong cửa sổ W ngày \\
outflow\_\allowbreak \{W\}d & Hành vi giao dịch & GIAO DỊCH & num & Tổng dòng tiền ra trong cửa sổ W ngày \\
net\_\allowbreak flow\_\allowbreak \{W\}d & Hành vi giao dịch & GIAO DỊCH & num & Dòng tiền thuần trong cửa sổ W ngày \\
avg\_\allowbreak txn\_\allowbreak amt\_\allowbreak \{W\}d & Hành vi giao dịch & GIAO DỊCH & num & Giá trị giao dịch trung bình trong cửa sổ W ngày \\
txn\_\allowbreak amt\_\allowbreak std\_\allowbreak \{W\}d & Hành vi giao dịch & GIAO DỊCH & num & Độ lệch chuẩn giá trị giao dịch trong cửa sổ W ngày \\
salary\_\allowbreak cnt\_\allowbreak \{W\}d & Hành vi giao dịch & GIAO DỊCH & num & Số lần nhận lương trong cửa sổ W ngày \\
salary\_\allowbreak amt\_\allowbreak \{W\}d & Hành vi giao dịch & GIAO DỊCH & num & Tổng lương nhận trong cửa sổ W ngày \\
repay\_\allowbreak amt\_\allowbreak \{W\}d & Hành vi giao dịch & GIAO DỊCH & num & Tổng tiền trả nợ trong cửa sổ W ngày \\
risky\_\allowbreak amt\_\allowbreak \{W\}d & Hành vi giao dịch & GIAO DỊCH & num & Chi tiêu nhóm rủi ro trong cửa sổ W ngày \\
cash\_\allowbreak wd\_\allowbreak cnt\_\allowbreak \{W\}d & Hành vi giao dịch & GIAO DỊCH & num & Số lần rút tiền mặt trong cửa sổ W ngày \\
cash\_\allowbreak wd\_\allowbreak amt\_\allowbreak \{W\}d & Hành vi giao dịch & GIAO DỊCH & num & Giá trị rút tiền mặt trong cửa sổ W ngày \\
ecom\_\allowbreak amt\_\allowbreak \{W\}d & Hành vi giao dịch & GIAO DỊCH & num & Chi tiêu thương mại điện tử trong cửa sổ W ngày \\
utility\_\allowbreak amt\_\allowbreak \{W\}d & Hành vi giao dịch & GIAO DỊCH & num & Chi tiêu hóa đơn tiện ích trong cửa sổ W ngày \\
night\_\allowbreak txn\_\allowbreak cnt\_\allowbreak \{W\}d & Hành vi giao dịch & GIAO DỊCH & num & Số giao dịch ban đêm (0h-5h) trong cửa sổ W ngày \\
large\_\allowbreak txn\_\allowbreak cnt\_\allowbreak \{W\}d & Hành vi giao dịch & GIAO DỊCH & num & Số giao dịch giá trị lớn (>50 triệu) trong cửa sổ W ngày \\
mobile\_\allowbreak txn\_\allowbreak cnt\_\allowbreak \{W\}d & Hành vi giao dịch & GIAO DỊCH & num & Số giao dịch qua mobile banking trong cửa sổ W ngày \\
outflow\_\allowbreak ratio\_\allowbreak \{W\}d & Hành vi giao dịch & GIAO DỊCH & num & Tỷ lệ dòng tiền ra trên dòng tiền vào trong cửa sổ W ngày \\
risky\_\allowbreak ratio\_\allowbreak \{W\}d & Hành vi giao dịch & GIAO DỊCH & num & Tỷ trọng chi tiêu rủi ro trong cửa sổ W ngày \\
cash\_\allowbreak ratio\_\allowbreak \{W\}d & Hành vi giao dịch & GIAO DỊCH & num & Tỷ trọng rút tiền mặt trong cửa sổ W ngày \\
night\_\allowbreak ratio\_\allowbreak \{W\}d & Hành vi giao dịch & GIAO DỊCH & num & Tỷ trọng giao dịch ban đêm trong cửa sổ W ngày \\
repay\_\allowbreak burden\_\allowbreak \{W\}d & Hành vi giao dịch & GIAO DỊCH & num & Gánh nặng trả nợ trên dòng tiền vào trong cửa sổ W ngày \\
txn\_\allowbreak amt\_\allowbreak cv\_\allowbreak \{W\}d & Hành vi giao dịch & GIAO DỊCH & num & Hệ số biến thiên giá trị giao dịch trong cửa sổ W ngày \\
days\_\allowbreak since\_\allowbreak last\_\allowbreak txn & Hành vi giao dịch & GIAO DỊCH & num & Số ngày từ giao dịch gần nhất \\
txn\_\allowbreak cnt\_\allowbreak trend\_\allowbreak 30\_\allowbreak 90 & Hành vi giao dịch & GIAO DỊCH & num & Xu thế số giao dịch: 30 ngày so với mức trung bình 90 ngày \\
inflow\_\allowbreak trend\_\allowbreak 30\_\allowbreak 90 & Hành vi giao dịch & GIAO DỊCH & num & Xu thế dòng tiền vào: 30 ngày so với mức trung bình 90 ngày \\
inflow\_\allowbreak trend\_\allowbreak 90\_\allowbreak 180 & Hành vi giao dịch & GIAO DỊCH & num & Xu thế dòng tiền vào: 90 ngày so với mức trung bình 180 ngày \\
income\_\allowbreak verification\_\allowbreak ratio & Hành vi giao dịch & GIAO DỊCH & num & Tỷ lệ dòng tiền vào thực tế so với thu nhập khai báo \\
salary\_\allowbreak regularity & Hành vi giao dịch & GIAO DỊCH & num & Mức độ đều đặn của việc nhận lương trong 180 ngày \\
income\_\allowbreak stability & Hành vi giao dịch & GIAO DỊCH & num & Độ ổn định thu nhập (1 - hệ số biến thiên dòng tiền vào theo tháng) \\
dti & Biến tổ hợp & TÍNH TOÁN & num & Debt-\allowbreak to-\allowbreak Income: (trả góp mới + nghĩa vụ hiện hữu) / thu nhập \\
pti & Biến tổ hợp & TÍNH TOÁN & num & Payment-\allowbreak to-\allowbreak Income: trả góp khoản vay mới / thu nhập \\
amount\_\allowbreak to\_\allowbreak income & Biến tổ hợp & TÍNH TOÁN & num & Số tiền vay / thu nhập hàng tháng \\
outstanding\_\allowbreak to\_\allowbreak income & Biến tổ hợp & TÍNH TOÁN & num & Tổng dư nợ hiện hữu / thu nhập hàng tháng \\
disposable\_\allowbreak income & Biến tổ hợp & TÍNH TOÁN & num & Thu nhập còn lại sau trả nợ và chi phí sinh hoạt tối thiểu \\
disposable\_\allowbreak ratio & Biến tổ hợp & TÍNH TOÁN & num & Tỷ lệ thu nhập còn lại trên thu nhập \\
income\_\allowbreak per\_\allowbreak dependent & Biến tổ hợp & TÍNH TOÁN & num & Thu nhập bình quân đầu người trong hộ \\
credit\_\allowbreak history\_\allowbreak ratio & Biến tổ hợp & TÍNH TOÁN & num & Độ dài lịch sử tín dụng so với tuổi trưởng thành \\
debt\_\allowbreak per\_\allowbreak institution & Biến tổ hợp & TÍNH TOÁN & num & Dư nợ bình quân trên mỗi tổ chức tín dụng \\
dpd\_\allowbreak x\_\allowbreak institutions & Biến tổ hợp & TÍNH TOÁN & num & Tương tác giữa số ngày quá hạn và số TCTD \\
high\_\allowbreak dti\_\allowbreak short\_\allowbreak tenure & Biến tổ hợp & TÍNH TOÁN & bin & Cờ rủi ro: DTI > 0,5 và thâm niên công việc < 12 tháng \\
balance\_\allowbreak to\_\allowbreak installment & Biến tổ hợp & TÍNH TOÁN & num & Số dư bình quân / khoản trả góp hàng tháng \\
inflow\_\allowbreak to\_\allowbreak installment & Biến tổ hợp & TÍNH TOÁN & num & Dòng tiền vào 30 ngày / khoản trả góp hàng tháng \\
\end{longtable}

{\footnotesize\itshape Kiểu dữ liệu: num -- định lượng; cat -- định tính; bin -- nhị phân. Mỗi đặc trưng có hậu tố \code{\_\{W\}d} được tính trên 4 cửa sổ thời gian $W \in \{7, 30, 90, 180\}$ ngày; 23 dòng dạng này tương ứng 92 đặc trưng, đưa tổng số đặc trưng lên 154.}
\endgroup

